
\documentclass[10pt]{article}
\usepackage[utf8]{inputenc}
\usepackage[T1]{fontenc}
\usepackage{amsmath, amssymb, xcolor, geometry, enumitem, multicol, tcolorbox}

\definecolor{mainblue}{RGB}{0, 102, 204}
\geometry{a4paper, margin=1.2cm, top=1cm, bottom=3cm, footskip=1.2cm}

\begin{document}
\enlargethispage{2.5cm}

% Modern Header
\begin{tcolorbox}[colback=mainblue!5, colframe=mainblue, sharp corners, boxrule=0.5pt]
    \small \textbf{Unit:} Unit 0: Foundations of Algebra \hfill \textbf{Name:} \rule{3cm}{0.4pt} \\
    \small \textbf{Topic:} Rational Numbers & Decimals \hfill \textbf{Date:} \rule{2cm}{0.4pt} \\
    \centering \textbf{\large Foundation (Jalapeno)}
\end{tcolorbox}

\vspace{0.2cm}

% MCQ Section
\begin{multicols}{2}
\begin{enumerate}[label=\textbf{Q\arabic*.}, leftmargin=*, itemsep=6pt, parsep=0pt]

  \item \small What is the decimal equivalent of the fraction $\frac{3}{4}$?
  \begin{enumerate}[label=(\Alph*), leftmargin=0.6cm, nosep, itemsep=2pt]
    \item 0.25
    \item 0.5
    \item 0.75
    \item 0.8
    \item 0.9
  \end{enumerate}

  \item \small A basketball player scored $\frac{7}{10}$ of the team's points. What percent of the points did the player score?
  \begin{enumerate}[label=(\Alph*), leftmargin=0.6cm, nosep, itemsep=2pt]
    \item 50%
    \item 60%
    \item 65%
    \item 70%
    \item 75%
  \end{enumerate}

  \item \small Which of the following rational numbers is greatest: $\frac{2}{3}$, $\frac{3}{4}$, $\frac{5}{6}$?
  \begin{enumerate}[label=(\Alph*), leftmargin=0.6cm, nosep, itemsep=2pt]
    \item $\frac{2}{3}$
    \item $\frac{3}{4}$
    \item $\frac{5}{6}$
    \item $\frac{2}{3}$ and $\frac{3}{4}$ are equal
    \item None of the above
  \end{enumerate}

  \item \small A cyclist traveled $2.5$ miles in $\frac{1}{2}$ hour. How many miles did the cyclist travel per hour?
  \begin{enumerate}[label=(\Alph*), leftmargin=0.6cm, nosep, itemsep=2pt]
    \item 4 miles
    \item 5 miles
    \item 6 miles
    \item 7 miles
    \item 8 miles
  \end{enumerate}

  \item \small What is the fraction equivalent of the decimal $0.6$?
  \begin{enumerate}[label=(\Alph*), leftmargin=0.6cm, nosep, itemsep=2pt]
    \item $\frac{1}{10}$
    \item $\frac{1}{5}$
    \item $\frac{3}{5}$
    \item $\frac{2}{3}$
    \item $\frac{3}{4}$
  \end{enumerate}

  \item \small A tennis player won $\frac{4}{5}$ of her matches. What percent of her matches did she win?
  \begin{enumerate}[label=(\Alph*), leftmargin=0.6cm, nosep, itemsep=2pt]
    \item 60%
    \item 70%
    \item 75%
    \item 80%
    \item 85%
  \end{enumerate}

  \item \small Which of the following rational numbers is least: $\frac{1}{2}$, $\frac{1}{3}$, $\frac{2}{5}$?
  \begin{enumerate}[label=(\Alph*), leftmargin=0.6cm, nosep, itemsep=2pt]
    \item $\frac{1}{2}$
    \item $\frac{1}{3}$
    \item $\frac{2}{5}$
    \item $\frac{1}{2}$ and $\frac{1}{3}$ are equal
    \item None of the above
  \end{enumerate}

  \item \small A swimmer finished a lap in $1.8$ minutes. What fraction of an hour did the swimmer take to finish the lap?
  \begin{enumerate}[label=(\Alph*), leftmargin=0.6cm, nosep, itemsep=2pt]
    \item $\frac{1}{20}$
    \item $\frac{1}{30}$
    \item $\frac{1}{50}$
    \item $\frac{3}{100}$
    \item $\frac{1}{100}$
  \end{enumerate}

\end{enumerate}
\end{multicols}

\vspace{-0.2cm}
\hrule
\vspace{0.2cm}
\textbf{\small Open Ended Questions (FRQ)}
\begin{enumerate}[label=\textbf{Q\arabic*.}, start=9, itemsep=5pt, topsep=0pt]

  \item \small Represent the rational numbers $\frac{1}{4}$, $\frac{1}{2}$, and $\frac{3}{4}$ on a number line and explain your reasoning. \vspace{2cm}

  \item \small A gymnast scored $85$ points out of $100$. Express this score as a fraction and a decimal. Calculate the percentage score and show your work. \vspace{2cm}

\end{enumerate}

\vfill

% Chef's Tips Footer
\begin{tcolorbox}[colback=blue!5, colframe=mainblue!50, title=\small \textbf{Chef's Tips}, fonttitle=\bfseries, bottom=1mm]
\begin{itemize}[noitemsep, topsep=2pt, leftmargin=0.5cm]
  \item \scriptsize Ensure students understand the concept of rational numbers and their representation on a number line.\item \scriptsize Encourage students to show their work and explain their reasoning for FRQs.\item \scriptsize Remind students to convert between fractions, decimals, and percents accurately.
\end{itemize}
\end{tcolorbox}

\end{document}
  